\section{Workflows}

Les quatre premiers workflows décrivent l'expérience de l'utilisateur ; ils sont identiques pour Blender et pour le web, seuls les contenus des observations et des actions changent. Les trois derniers décrivent notre façon de travailler.

% Notes pgf-umlsd : un libellé de \mess contenant des parenthèses doit être entre doubles accolades ;
% un tiret bas dans un libellé de \mess doit passer par \detokenize{...} ; pas de caractères actifs
% de babel français (: ; ! ? et guillemets) dans les libellés.

\subsection{Expliquer ce que l'utilisateur regarde}

Exemple : dans Blender, l'utilisateur sélectionne la citrouille et demande « Explique-moi ces modificateurs ».

\begin{center}
\resizebox{0.97\textwidth}{!}{%
\begin{sequencediagram}
\newthread{u}{Utilisateur}
\newinst[1.2]{p}{Panneau}
\newinst[1.6]{c}{Cœur}
\newinst[1.8]{l}{LLM}
\newinst[1.8]{a}{Adaptateur}
\mess{u}{{question (voix ou texte)}}{p}
\mess{p}{user.message}{c}
\begin{call}{c}{get\_context, observe(focus)}{a}{contexte, observation}
\end{call}
\begin{call}{c}{contexte + outils}{l}{show\_entities(refs, stack)}
\end{call}
\begin{callself}{c}{projection, validateur}{}
\end{callself}
\mess{c}{{board.patch (nœuds observed)}}{p}
\begin{call}{c}{résultat de l'outil}{l}{narrate(segments, repères)}
\end{call}
\mess{c}{narration.segment}{p}
\mess{c}{{adapter.highlight(ref)}}{a}
\mess{p}{voix, sous-titres, focus}{u}
\end{sequencediagram}}
\end{center}

Points clés : le LLM choisit \emph{quelles} entités montrer, la projection fabrique les nœuds à partir de l'observation réelle. Un concept général (par exemple « ordre d'évaluation ») est ajouté avec \code{add\_concept} en provenance \pinf{} ou \pill. Les repères des segments synchronisent la surbrillance sur le tableau et dans Blender avec la phrase prononcée.

\subsection{Action proposée, permission, exécution et vérification}

Exemple : « Arrondis les arêtes de la bougie. » Le LLM propose \code{add\_modifier(Candle, BEVEL)} en mode \code{do}.

\begin{center}
\resizebox{0.97\textwidth}{!}{%
\begin{sequencediagram}
\newthread{u}{Utilisateur}
\newinst[1.2]{p}{Panneau}
\newinst[1.6]{c}{Cœur et porte}
\newinst[1.8]{l}{LLM}
\newinst[1.8]{a}{Adaptateur}
\mess{l}{{\detokenize{propose_action(add_modifier, mode do)}}}{c}
\begin{callself}{c}{valider paramètres, politique}{risque reversible}
\end{callself}
\mess{c}{permission.request}{p}
\mess{p}{dialogue Approuver ou Refuser}{u}
\mess{u}{clic ou oui vocal}{p}
\mess{p}{permission.respond}{c}
\begin{call}{c}{adapter.execute}{a}{ok, undo\_token}
\end{call}
\mess{a}{{\detokenize{event.state_changed (origin wysp)}}}{c}
\begin{call}{c}{adapter.observe}{a}{observation fraîche}
\end{call}
\begin{callself}{c}{post-conditions}{preuve}
\end{callself}
\mess{c}{board.patch, evidence.added}{p}
\mess{c}{résultat vérifié}{l}
\end{sequencediagram}}
\end{center}

\begin{itemize}[leftmargin=1.2em]
  \item \textbf{Échec de vérification} : le cœur appelle \code{adapter.undo} avec le jeton, puis explique l'échec avec la valeur attendue et la valeur constatée.
  \item \textbf{Mode guidé} (préféré pour apprendre) : aucun \code{execute}. Le cœur publie les étapes de \code{guide.ui\_path} (« Propriétés › Modificateurs › Add Modifier › Bevel »), met en évidence chaque cible, et vérifie les post-conditions dès que l'adaptateur signale un changement d'origine utilisateur.
  \item \textbf{Annuler} : « annule ça » rejoue le dernier jeton de la pile d'annulation.
\end{itemize}

\subsection{Interruption et reprise}

\begin{enumerate}[leftmargin=1.2em]
  \item Pendant la narration, l'utilisateur parle (détection d'énergie) ou appuie sur Échap. Le panneau coupe l'audio et envoie \code{user.interrupt} (moins de 150\,ms).
  \item Le cœur annule le tour (jeton d'annulation vérifié entre chaque étape et pendant le flux du LLM), publie \code{narration.cancel} et garde le curseur : dernier segment entendu. Une action déjà commencée chez l'adaptateur n'est jamais coupée en deux ; sa vérification est enregistrée.
  \item Si l'utilisateur pose une question (« Pourquoi Solidify après Subdivision ? »), un nouveau tour démarre avec le nœud concerné dans le contexte. S'il dit « continue », le cœur republie exactement les segments non entendus, sans nouvel appel au LLM.
  \item Le bouton « Pourquoi ? » d'un nœud produit le même effet qu'une question sur ce nœud.
\end{enumerate}

\subsection{Exercice et preuve}

\begin{center}
\resizebox{0.97\textwidth}{!}{%
\begin{sequencediagram}
\newthread{u}{Utilisateur}
\newinst[1.2]{p}{Panneau}
\newinst[1.6]{e}{Moteur pédagogique}
\newinst[1.6]{v}{Vérification}
\newinst[1.6]{a}{Adaptateur}
\begin{callself}{e}{liaisons résolues, critère encore faux}{offre}
\end{callself}
\mess{e}{{exercise.update (consigne)}}{p}
\mess{u}{modifie la scène lui-même}{a}
\mess{a}{{\detokenize{event.state_changed (origin user)}}}{e}
\begin{call}{e}{vérifier les prédicats}{v}{preuve}
\begin{call}{v}{adapter.observe}{a}{observation}
\end{call}
\end{call}
\begin{callself}{e}{mise à jour de la maîtrise}{}
\end{callself}
\mess{e}{{evidence.added (verified), learner.update}}{p}
\mess{p}{vérifié dans l'application}{u}
\end{sequencediagram}}
\end{center}

Sans succès, l'utilisateur peut demander un indice (révélés dans l'ordre). Après \code{max\_attempts} échecs, Wysp propose la solution en mode guidé ou \code{do}. Si l'application ne permet pas de vérifier, la déclaration de l'utilisateur est enregistrée comme preuve \code{declared}, affichée comme non vérifiée.

\subsection{Développer avec Kiro}

\begin{figure}[h]
\centering
\begin{tikzpicture}[font=\sffamily\scriptsize, node distance=0.35cm and 0.45cm,
  st/.style={draw=#1, fill=#1!8, rounded corners=2pt, align=center, minimum height=0.95cm, text width=2.15cm},
  arr/.style={-{Stealth[length=2mm]}, thick, color=surface}]
  \node[st=pa] (spec) {Spec\\\tiny exigences, design, propriétés};
  \node[st=pa, right=of spec] (task) {Tâche\\\tiny « Start task »};
  \node[st=pumpkin, right=of task] (pre) {Hook\\\tiny task-preflight};
  \node[st=pb, right=of pre] (impl) {Implémentation\\\tiny agent + humain};
  \node[st=pumpkin, right=of impl] (save) {Hooks à la sauvegarde\\\tiny format, frontière, schémas, contrat};
  \node[st=pumpkin, below=0.9cm of save] (post) {Hook\\\tiny task-tests};
  \node[st=pc, left=of post] (guard) {Agent\\\tiny contract-guardian};
  \node[st=pb, left=of guard] (pr) {Pull request\\\tiny revue du responsable};
  \node[st=okgreen, left=of pr] (ci) {CI verte\\\tiny tests, kit, frontière};
  \node[st=okgreen, left=of ci] (merge) {Fusion\\\tiny au moins 1 fois par jour};
  \draw[arr] (spec) -- (task); \draw[arr] (task) -- (pre); \draw[arr] (pre) -- (impl); \draw[arr] (impl) -- (save);
  \draw[arr] (save) -- (post); \draw[arr] (post) -- (guard); \draw[arr] (guard) -- (pr); \draw[arr] (pr) -- (ci); \draw[arr] (ci) -- (merge);
  \draw[arr, dashed, color=muted] (save.north) |- ++(0,0.45) -| node[pos=0.25, above, text=muted] {échec : correction immédiate} (impl.north);
\end{tikzpicture}
\caption{Boucle de développement d'une tâche de spec.}
\end{figure}

\begin{enumerate}[leftmargin=1.2em]
  \item Chaque personne ouvre la spec de son périmètre dans Kiro et lance les tâches dans l'ordre (« Start task »). Le hook \code{task-preflight} rappelle les exigences, le responsable et le processus de changement de contrat.
  \item Pendant l'implémentation, les hooks de sauvegarde formatent le code et lancent les contrôles déterministes. Ils ne coûtent rien quand rien n'est concerné.
  \item À la fin de la tâche, \code{task-tests} exécute les tests des paquets modifiés. Les tâches de propriétés sont obligatoires.
  \item Avant la pull request, l'agent \code{contract-guardian} (lecture seule) produit un tableau de conformité. Le responsable du périmètre touché relit. La CI doit être verte.
  \item Le journal \file{docs/kiro-journal.md} note ce que Kiro a produit et ce qui a été corrigé à la main : c'est la matière de la section « How we used Kiro » de la soumission.
\end{enumerate}

\subsection{Changer le contrat}

\begin{enumerate}[leftmargin=1.2em]
  \item Écrire le besoin dans la spec concernée, puis modifier \file{.kiro/specs/adapter-contract/} (exigence et design).
  \item Modifier les modèles Pydantic dans \file{core/src/wysp_core/contract/} et augmenter \code{contract\_version} (mineure si ajout optionnel).
  \item Le hook \code{contract-schema-sync} régénère \file{schemas/} et les types TypeScript.
  \item Lancer \code{wysp-contract --offline} pour les deux tuteurs : les enregistrements montrent qui serait cassé.
  \item Obtenir l'accord de \PA{} et du responsable de chaque adaptateur touché. Après J0, pas de changement majeur sans accord des trois.
\end{enumerate}

\subsection{Créer un nouveau tuteur avec Wysp}

C'est le parcours que le jury doit voir comme valeur du squelette, et que le README présentera en quatre étapes.

\begin{enumerate}[leftmargin=1.2em]
  \item \textbf{Manifeste} : identifiant, niveau d'intégration, capacités, types d'entités, schéma de références.
  \item \textbf{Observation} : traduire l'état de l'application en entités (le SDK fournit client, empreinte et validation).
  \item \textbf{Actions} : un fichier par action avec schéma, risque, post-conditions et chemin d'interface ; le hook \code{new-action-integration} câble le reste.
  \item \textbf{Profil} : compétences, concepts, leçons et exercices en YAML, validés par \code{wysp validate-profile}.
\end{enumerate}
Puis \code{wysp-contract --tutor <nom>} indique si l'adaptateur est conforme. Aucun changement du cœur ni du panneau n'est nécessaire.
